% Modernised reading — fonds Alexandre Grothendieck, shelfmark 37, pages 1-99.
%
% This is an INTERPRETATION, not a transcription. It re-reads the pages in
% today's notation and vocabulary, states the hypotheses the manuscript leaves
% implicit, and drops the critical apparatus so that the mathematics reads
% straight through. Everything the brackets and underlines carried moves into
% a footnote. Where the manuscript is loose or mistaken, the true statement is
% what appears here and the footnote says what the page has. In any dispute of
% substance the transcription governs: whoever cites Grothendieck must cite the
% batch-NN.fr.tex files.
%
% Pass: Opus 5.5 (claude-opus-5-5), 2026-10-04 — first pass, unchecked against
% the pages by a human. The folder was read whole: five batches, pages 1-99,
% all transcribed. This reading was made on Opus 5.5 and is not measured
% against the Opus 5 reference reading of folder 115.
%
% WHAT THE FOLDER IS. A dossier of editorial work, not a mathematical text:
% the Springer re-editions of SGA 1, 4, 5, 6, 7 and of EGA I and II, ordered
% by tan covers (« SGA 1 », « SGA 4 », « SGA 5 », « SGA 6 », « SGA 7 »,
% « Réédition EGA I », « EGA II »), with letters, lists of corrections and a
% few pieces of unpublished mathematics. The spine is therefore the order of
% the covers, and the mathematics is read where it sits on it.
%
% THIRD-PARTY LETTERS AND TEXTS (Verdier pp. 8-11, 13; Artin pp. 33-34;
% Hartshorne pp. 53-54; an unidentified reader's commentary pp. 65-71;
% Deligne pp. 74-75; Dieudonné p. 95) are noted only in the transcription,
% and this reading restates nothing of them beyond what those notes say
% (RIGHTS.md). His own letters and typescripts are read in full.
%
% DATING DOUBTS kept visible (src/content/letters.json): Verdier's « 16
% Juillet » (p. 8) and « le 28 Mai » (p. 13) carry no year, 1971 is inferred;
% p. 12 thanks for a letter « du 23 Mai », to be settled on the facsimile;
% the letter on « 2.1.1 » (pp. 29-30) and Artin's (pp. 33-34) are undated,
% their placing about 1969-1971 is not established; pp. 2, 17 and all the
% manuscript notes are undated.
%
% NOTATION fixed once: \widetilde{X}_{\text{ét}} for the étale topos;
% \widehat{C} for presheaves on C (the pages write C^\wedge); Ind(C); Ouv(E)
% for subobjects of the final object; L for a set of primes; A_n = A/T^{n+1}
% in the EGA exchange. Exposé numbers are roman, « SGA 4 VI 1.25 ».
%
% Substantive departures from the pages, each footnoted where it occurs:
%   - p. 16: hypothesis (ii) as written is weaker than the comparison lemma's,
%     and the page announces, does not prove, the cocontinuity;
%   - p. 22: (i bis) says « adjoint à droite »; the condition that makes
%     Ind(C) a topos is a LEFT EXACT LEFT adjoint — read so;
%   - pp. 29-30: L must consist of primes invertible on the base (on k in
%     2.1.2) — supplied;
%   - p. 37: the Čech-type formula by sieves is false in general for n >= 2;
%     hypercoverings (Verdier) give the true one;
%   - p. 48: « all morphisms of C are monos » is false as a criterion
%     (C = a group); the criterion is « C is a preorder »; the first line's
%     illegible word is not guessed;
%   - p. 51: D(i_* F) = i_* D(F) holds for the underived direct image in
%     dimension 1, not in general — supplied;
%   - p. 92: the written generators of J_n are NOT linearly independent over
%     B (X_0 . X_1 T lies in T J_n); they are independent modulo the ideal
%     (X_i), which is what the minimality needs;
%   - p. 96: the lemma's blank is \sum m_i; A_{f_i} should be the completed
%     ring A_{\{f_i\}}; « tensorisant par A_0 » is by A_n.
%
% Announced and never established in the folder: the 5°) of exercise 4.10.6
% (p. 16 breaks off); exercise « 2 » on Ind(C) (p. 22) and the exercises of
% VIII § 10 (pp. 38-39) are stated, not proved; 10.5 is empty; the letter of
% p. 93 breaks off mid-sentence; whether a) => c) in EGA I 10.10.5 holds
% without noetherian hypotheses is left open by both correspondents.

\documentclass[11pt,a4paper]{article}
\input{../preamble/grothendieck}

\folder{37}
\pages{1}{99}
\dating{1967-1971}
\watermark{Édition de démonstration\\interprétation personnelle de l'œuvre}
\foldertitle{Rééditions SGA, EGA [SGA 4 à SGA 7, EGA I et EGA II] : tapuscrits et copies de tapuscrit annoté (s.d.), lettres (1967-1971, s.d.), notes manuscrites (s.d.) — lecture modernisée du dossier entier}

\begin{document}

\section*{Résumé}

\begin{resume}
Ce dossier montre l'envers d'une œuvre publiée : comment des séminaires de
recherche, tirés à l'origine en quelques centaines d'exemplaires ronéotés à
l'Institut des Hautes Études Scientifiques, deviennent des livres. Entre 1967
et 1971, les \emph{Séminaires de Géométrie Algébrique} (SGA) et les
\emph{Éléments de Géométrie Algébrique} (EGA) sont réédités chez Springer.
On voit ici ce que cela demande : retrouver les exemplaires d'archives,
relancer des rédacteurs dispersés (Verdier à Brandeis, Ferrand qui ne rend pas
son exposé), compter les pages, faire taper, relire, corriger des renvois
ligne par ligne, et décider de ce qu'on laisse incomplet.

Le moment est particulier. En septembre 1970 l'auteur quitte l'IHÉS, parce
que l'institut accepte des subventions militaires ; la dernière page d'une
préface, dans un état antérieur au texte imprimé, le dit sans détour. Huit mois
plus tard il écrit qu'il ne trouve « pas une minute » pour la réédition de
SGA 4, et qu'il est prêt à la laisser « incomplète et dégueulasse » plutôt que
de la payer d'un jour ou deux pris sur le « boulot survie, infiniment plus
important ». Le dossier enregistre ainsi, en marge des mathématiques, le
moment où celui qui les avait organisées commence à s'en détourner --- et la
conscience professionnelle avec laquelle il continue pourtant à poser des
questions précises sur des exercices faux.

Car il y a des mathématiques, et quelques-unes ne sont pas dans les livres
publiés. L'idée qui les relie toutes est celle de \emph{topos} : au lieu
d'étudier un espace par ses points, on l'étudie par la catégorie des
« faisceaux » qu'il porte --- de la même façon qu'on peut connaître un groupe
par toutes ses représentations plutôt que par ses éléments. Cette catégorie
garde l'essentiel de l'espace, et elle existe dans des situations où les
points font défaut. Quatre morceaux s'en détachent.

D'abord un dictionnaire, fait à la main sur le verso d'une page de SGA 4,
entre quatre types d'objets finis : les espaces topologiques finis, les
ensembles ordonnés finis, les treillis finis, et les topos « finis ». C'est
le cas le plus élémentaire où l'on puisse voir qu'un espace, l'ordre de ses
points et la forme de ses ouverts sont trois façons de dire la même chose.

Ensuite, un paragraphe resté inédit de SGA 4 qui définit les \emph{étendues
algébriques} et les \emph{espaces algébriques} comme des topos annelés
localement semblables au topos étale d'un schéma --- l'idée, aujourd'hui
centrale, qu'un objet géométrique peut être recollé non par des ouverts mais
par des « revêtements » étales, quitte à garder la trace de ses symétries.

Puis un exercice, inédit lui aussi, qui demande quand la catégorie des
« ind-objets » d'une petite catégorie est un topos, et une lettre qui
propose d'énoncer pour des \emph{champs} (des faisceaux à valeurs dans les
catégories) le théorème de changement de base lisse.

Enfin, en 1967, un échange de lettres avec Dieudonné sur la façon d'énoncer,
sans hypothèse noethérienne, un critère des EGA : Grothendieck y construit un
contre-exemple, prouve un lemme sur le nombre de générateurs d'un module, et
signale qu'une démonstration d'EGA IV est en défaut, la question se ramenant à
savoir si un module « localement projectif » est projectif. Ni l'un ni
l'autre ne tranche ; la réponse viendra de Raynaud et Gruson.

Ce que cherche un lecteur qui veut aller plus loin : la théorie des topos
(topos cohérents, théorème de Deligne sur les points), la topologie
d'Alexandrov et la dualité de Birkhoff pour le dictionnaire fini, les espaces
algébriques et les champs de Deligne--Mumford pour les étendues, et la
descente de la projectivité pour la dernière question.

\keywords{SGA re-edition, étale topos, coherent topos, locally finitely presentable topos, Ind-completion, smooth base change, Giraud stacks, hypercovering, Alexandrov topology, Birkhoff duality, algebraic space, Deligne-Mumford stack, étendue, formal scheme, formally smooth morphism, descent of projectivity, cotangent complex, local invariant cycle theorem}
\end{resume}

\section*{Le fil du dossier, et ses conventions}

Le dossier est rangé par séminaire, sous des chemises de papier bistre
inscrites en haut à droite ; ces chemises donnent son fil, et la lecture le
suit :

\begin{itemize}
\item pages 2 à 4 --- la réédition dans son ensemble : ce qui reste à
frapper, deux lettres du 3.11.1970 ;
\item « SGA 1 », page 6 --- la dernière page d'une préface ;
\item « SGA 4 », pages 8 à 49 --- de loin la plus longue section :
l'exercice IV 4.10.6, le sommaire de VI, le plan en trois fascicules, les
notes exposé par exposé, un exercice inédit sur $\mathrm{Ind}(C)$, la lettre
sur « 2.1.1 », la réponse d'Artin, un § 10 inédit de VIII sur les étendues
algébriques, la liste de terminologie, le dictionnaire des topos finis ;
\item « SGA 5 », pages 51 à 54 --- notes, et la lettre de Hartshorne ;
\item « SGA 6 », pages 56 à 71 --- notes, et les commentaires dactylographiés
d'un lecteur sur l'exposé X ;
\item « SGA 7 », pages 73 à 82 --- la lettre à Ferrand (qui concerne en fait
SGA 6), la lettre de Deligne de 1968, les appendices prévus ;
\item « Réédition EGA I », pages 84 à 97 --- la table des matières de la
nouvelle édition, et l'échange avec Dieudonné d'août-septembre 1967 ;
\item « EGA II », page 99.
\end{itemize}

Les dates écrites vont du 27.8.1967 au 23.6.1971 ; elles ne suivent pas
l'ordre des pages. Plusieurs pièces n'ont \emph{pas d'année} et leur datation
est inférée : les lettres de Verdier « 16 Juillet » (page 8) et « le 28 Mai »
(page 13), données pour 1971 parce qu'elles répondent à la lettre du
18.5.1971 ; la lettre sur « 2.1.1 » (pages 29-30) et celle d'Artin (pages
33-34), sans date du tout. Ces doutes sont rappelés là où ils comptent.

Les lettres et textes d'autres personnes --- Verdier, Artin, Hartshorne,
Deligne, Dieudonné, un lecteur anonyme de SGA 6 --- ne sont pas reproduits
dans la transcription, qui en donne seulement une notice ; cette lecture
n'en dit pas plus que ces notices. Ses lettres à lui sont lues en entier.

Notations : $\widetilde{X}_{\text{ét}}$ est le topos étale du schéma $X$ ;
$\widehat{C}$ la catégorie des préfaisceaux d'ensembles sur une petite
catégorie $C$ (les pages écrivent $C^{\wedge}$) ; $\mathrm{Ouv}(E)$ l'ensemble
ordonné des sous-objets de l'objet final d'un topos $E$ ; $\mathrm{Ind}(C)$ la
catégorie des ind-objets de $C$. Un renvoi « VI 1.25 » se lit dans le
séminaire de la section où l'on se trouve.

Ce que le dossier annonce sans l'établir : la suite (5°) de l'exercice
4.10.6, interrompue ; les démonstrations de l'exercice sur $\mathrm{Ind}(C)$
et des exercices du § 10 de VIII, seulement énoncés ; et, dans l'échange avec
Dieudonné, l'implication a) $\Rightarrow$ c) de EGA I 10.10.5 sans hypothèse
noethérienne, que ni l'un ni l'autre n'arrive à prouver ou à réfuter.

\pagerange{2}{4}
\section*{La réédition dans son ensemble (pages 2 à 4)}

Une feuille dactylographiée, sans date ni signature, intitulée « SGA :
Manuscripts à frapper », fait l'état des lieux séminaire par séminaire. Pour
SGA 4 : la fin de IV, V et la fin de VI par Verdier, VI B et un appendice à
XVII par Saint-Donat, les tables et l'index ; pour SGA 5 : IV (rédacteur
inconnu), VII par Jouanolou, XI et XII par Bucur, XIV par lui-même ; pour
SGA 6 : XI par Ferrand ; pour SGA 7 : des exposés sans rédacteur désigné, V
par Mme Raynaud, VI par Rim. Les nombres de pages sont estimés ; un nombre
souligné signale un manuscrit disponible, un point d'interrogation un exposé
qui ne sera peut-être jamais rédigé.\footnote{Au crayon, vraisemblablement de
sa main sans que ce soit établi : la ligne de SGA 1 (un index) est barrée,
suivie de deux mots dont le dernier se lit peut-être « Saarbach » ; « écrire
Bucur » en face de SGA 5 ; un trait ondulé qui annule tout le bloc de SGA 6 ;
« écr. Deligne » en face de SGA 7 V et VI.}
Le total est d'environ 250 pages prêtes, et 200 à venir pour SGA 4,
« cela dépendra de Mr. Verdier, qui est à Brandeis cette année ».\footnote{La
feuille n'est pas datée ; la mention de Verdier à Brandeis et l'accord avec
les lettres du 3.11.1970 la placent à l'automne 1970, sans que ce soit
établi.}

Deux doubles de lettres du 3.11.1970 suivent. À Mme Michon, à l'IHÉS : il
manque, parmi les exemplaires d'archives qu'il a emportés, SGA 4 XVII et la
première partie de VI, dont il a besoin pour « l'édition photooffset en
préparation ». À Deligne : a-t-il fini de lire l'exposé de Rim (SGA 7 VI) ?
corrigera-t-il la frappe de son exposé SGA 4 XVIII ? a-t-il des nouvelles de
Ferrand pour SGA 6 XI ?

\pagerange{6}{6}
\section*{SGA 1 : la fin d'une préface (page 6)}

Sous la chemise « SGA 1 », un seul feuillet tapé, paginé (ix), signé « Massy,
Août 1970 » : la dernière page de la préface de la réédition de SGA 1, dans
un état antérieur au texte imprimé.\footnote{La transcription l'a comparé à
la préface publiée (Lecture Notes 224) : celle-ci dit « immoral et
dangereux », non « profondément immoral et extrêmement dangereux » ; « quitte
l'IHES le 30 septembre 1970 », non « le 30 Septembre prochain » ; elle
attribue la décision de ne plus diffuser à une demande de l'auteur
(« J'ai demandé à M. Motchane \ldots »), là où ce feuillet dit que le
directeur « s'est déclaré d'accord » ; et elle ajoute un remerciement à
J.~P. Delale pour les index.}
Il y explique s'être trouvé travailler trois ans dans une institution « alors
qu'elle participait à [son] insu à un mode de financement » qu'il juge
« profondément immoral et extrêmement dangereux » ; seul de cet avis parmi
ses collègues, il quitte l'IHÉS et suspend toute collaboration scientifique
avec lui tant qu'il acceptera de telles subventions ; l'IHÉS cessera de
diffuser ses textes et le Séminaire, que Springer publiera dans les
\emph{Lecture Notes}, en prenant à sa charge la frappe des exposés nouveaux
et manquants. C'est ce qui donne au dossier entier sa raison d'être.

\pagerange{8}{16}
\section*{SGA 4 : l'exercice IV 4.10.6 (pages 8 à 16)}

\subsection*{Une correspondance autour d'un exercice faux}

Le premier objet de la section SGA 4 est un exercice. La lettre du 3.11.1970
à Verdier (page 15) rappelle ce qui manque pour que SGA 4 soit complet : V, et
VI à partir du § 5, que Verdier rédige ; VI B de Saint-Donat, qu'il doit
relire ; l'appendice à XVII, que Grothendieck lui envoie. Six mois plus tard,
le 18.5.1971 (page 14), la réédition est « bloquée » : il manque les deux
appendices de Deligne (« assez de points », « image inverse d'un plat ») et
une version corrigée de l'exercice IV 4.10.6, dont il a « entièrement
oublié ce qui était canulé ». « C'est un exercice qui sert tout le temps
(!), il ne devrait donc pas être question de le vider. Il doit y avoir
quelque chose de vrai à mettre à la place du faux \ldots » Le même jour il
demande à Verdier de demander à Artin ce qu'il faut faire des exposés
postérieurs à VI.

Verdier répond (page 13), d'une lettre que la transcription ne reproduit
pas : selon sa notice, les deux appendices iront aux exposés V et VI, une
« version non canulée de 4.10.6 » suivra, et il demande une copie de IV
§§ 9-13.\footnote{Cette lettre est datée « le 28 Mai » sans année ; 1971 se
déduit de l'en-tête de Brandeis et de la lettre à laquelle elle répond. La
réponse de Grothendieck du 23.6.1971 (page 12) remercie pour une lettre
« du 23 Mai » --- la page porte bien 23. Laquelle des deux dates est la
bonne reste à vérifier sur le fac-similé.}
Le 23.6.1971 (page 12), l'exercice n'est toujours pas arrivé : la dactylo
termine, Springer rend la machine louée, et la publication du fascicule 1 de
SGA 4 est remise « sine die ». La nouvelle rédaction arrive enfin avec une
lettre de Verdier datée « 16 Juillet », sans année (page 8) ; d'après la
notice de la transcription, elle l'envoie avec des broutilles relevées dans
IV n°s 9 sqq., et signale qu'il manque le n° 14 ; Grothendieck note en marge
« demander à Messing ».\footnote{L'année 1971 est inférée, non écrite : la
lettre envoie l'exercice corrigé que réclament les lettres du 18.5.1971 et du
23.6.1971.}
Le tapuscrit joint (pages 9 à 11) est la rédaction de Verdier ; la
transcription n'en donne que l'objet --- un $\mathcal{U}$-site $S$ de
topologie moins fine que la canonique, une partie $\underline{M}$ de
$\mathrm{Fl}(S)$ soumise à des conditions a) à d), le « petit site »
$S(X)$, les foncteurs $\mathrm{Res}_X$ et $\mathrm{Prol}_X$, le plongement du
« petit topos » $S(X)^{\sim}$ dans le « gros topos » $(S_{/X})^{\sim}$, puis
l'exercice 4.10.7 ---, et il n'en est pas dit plus ici.\footnote{Le modèle
est celui du topos étale d'un schéma $X$ (petit topos, faisceaux sur les
$X$-schémas étales) plongé dans le topos de tous les $X$-schémas munis de la
topologie étale (gros topos). Ce rapprochement est le nôtre.}

\subsection*{Une généralisation, inachevée}

Sur papier de l'IHÉS, une page (16) intitulée « Continuation de SGA 4 IV
4.10.6 » commence une généralisation des points 3° et 4° de l'exercice :
au lieu du petit site $S(X) \subset S_{/X}$, une sous-catégorie pleine
quelconque. Lue en clair :

\emph{Soit $S_0 \subset S$ une sous-catégorie pleine telle que}

(i) \emph{si $u : X \to Y$ est dans $\underline{M}$ et $Y \in S_0$, alors
$X \in S_0$ ;}

(ii) \emph{pour tout $X \in S$ il existe une famille couvrante $(X_i \to X)$
dont chaque $X_i$ admet un morphisme vers un objet de $S_0$.}

On munit $S_0$ de la topologie induite. Le feuillet se propose de prouver
qu'une famille de morphismes de $S_0$ est couvrante dans $S_0$ si et
seulement si elle l'est dans $S$ --- ce qu'il paraphrase en disant que
l'inclusion $S_0 \to S$ est cocontinue ---, puis que l'ensemble
$\underline{M}_0$ des $u : X \to Y$ de $\underline{M}$ avec $Y \in S_0$ (donc
$X \in S_0$ par (i)) satisfait dans $S_0$ les conditions a) à d) de 1°, et
enfin de construire, comme au 3°, une restriction
$\mathrm{Res} : S^{\sim} \to S_0^{\sim}$ commutant aux limites inductives et
projectives et un prolongement
$\mathrm{Prol} : S_0^{\sim} \to S^{\sim}$ exact et pleinement fidèle.
Rien de cela n'est démontré : la feuille s'arrête au milieu de cette
dernière phrase.\footnote{La fin de la page est très cursive ; « cocontinu »,
« ssi », « commute », « pleinement fidèle » sont des lectures incertaines,
guidées par le 3° de l'exercice. Deux remarques de notre fait. L'hypothèse
(ii), telle qu'écrite (« le crible engendré par les objets de $S_0$ est
couvrant »), demande seulement que les $X_i$ \emph{s'envoient} dans des
objets de $S_0$ ; c'est plus faible que l'hypothèse du lemme de comparaison
(SGA 4 III 4.1 : tout objet de $S$ est couvert par des objets de $S_0$), qui
est ce qui donne d'ordinaire $S_0^{\sim} \simeq S^{\sim}$. Ici on ne veut
justement pas une équivalence, mais un plongement analogue à celui du petit
topos dans le gros, et c'est (i), la stabilité par les flèches de
$\underline{M}$, qui doit en tenir lieu ; le feuillet ne dit pas comment. Par
ailleurs, « couvrante dans $S_0$ si et seulement si couvrante dans $S$ » est
la propriété qui rend l'inclusion à la fois continue et cocontinue lorsque
$S_0$ est stable par les produits fibrés qu'on utilise ; elle n'est pas, à
elle seule, la définition de la cocontinuité.}

\pagerange{17}{17}
\section*{SGA 4 : le sommaire de l'exposé VI (page 17)}

Un sommaire dactylographié de l'exposé VI, « Conditions de finitude », repris
à l'encre, donne l'ordre qu'il voulait :

\begin{enumerate}
\item conditions de finitude pour les objets et flèches d'un topos (pages 1
à 19 du tapuscrit) ;
\item pour un topos (20 à 37) ;
\item pour un morphisme de topos (38 à 46) ;
\item pour un topos obtenu par recollement (47 à 65) ;
\item cohomologie d'une limite inductive de faisceaux sur un topos cohérent ;
\item catégories fibrées et limites inductives (rappels) ;
\item sites et topos fibrés ;
\item limites projectives de topos ;
\item cas des limites projectives filtrantes de topos cohérents à morphismes
de transition cohérents : passage à la limite en cohomologie ;
\item appendice, par P. Deligne : un topos cohérent a suffisamment de
points.\footnote{L'ordre ci-dessus est celui qui résulte des renumérotations
à l'encre : le n° 9 dactylographié devient 5 et remonte après le n° 4, les
n°s 5 à 7 deviennent 6 à 8 ; le numéro du « 8 » dactylographié est surchargé
d'un chiffre illisible, et nous le plaçons en 9 sous réserve. Le n° 4 tapé
était « Constructibilité » ; il est barré et remplacé à la main. La ligne 10
est entièrement manuscrite, d'une écriture très rapide ; « suffisamment de
points » est lu sous réserve. Que l'appendice soit le théorème de Deligne sur
les points des topos cohérents est une identification de notre fait,
cohérente avec la lettre du 18.5.1971 (« assez de points »).}
\end{enumerate}

Deux paragraphes dactylographiés, ensuite annulés au crayon, disaient que le
n° 4 contiendrait « le sorite des groupes constructibles », des faisceaux de
torsion, de la constructibilité pour les Modules et leurs complexes, et que
l'auteur pensait le rédiger lui-même « mais ne le garanti[t] pas ». Reste,
non barrée, l'annonce que la rédaction des n°s 1 à 3 est complète, à quelques
exemples près : un topos localement algébrique non algébrique, et --- ajouté
au crayon --- un topos ayant une famille génératrice d'objets cohérents qui
n'est pas localement algébrique.\footnote{Le feuillet est tenu pour un
tapuscrit de lui, inédit, d'après les arguments de la transcription : il
diffère du sommaire publié de VI, et son auteur y annonce rédiger lui-même un
numéro. Une note pâle au crayon en tête (« 1.9.5, 1.23, 2.9 --- à reporter au
n° 3 ») n'est lue qu'en partie.}

\pagerange{18}{20}
\section*{SGA 4 : renvois, fascicules, stencils (pages 18 à 20)}

Trois pages de travail. La première relève dans le tirage de SGA 4 XVII six
renvois à vérifier : à V (pages 75, 170, 175), à $f_{!}$ pour les ensembles
pointés et les faisceaux en groupes (« Consulter Verdier »), à VI B, et à la
seconde édition d'EGA II pour le lemme de Chow.

La deuxième est un plan de SGA 4 en trois fascicules, avec un nombre de pages
par exposé : le premier (préface, exposés I à IV, tables et index) d'environ
570 pages ; le second, exposés V à XI, de 432 pages (V : 100, VI : 150, VI B :
30, VII : 24, VIII : 46, IX : 42, X : 21, XI : 19 --- la somme est exacte) ;
le troisième, « Théorèmes de changement de base et de dualité globale »,
exposés XII à XIX, d'environ 640 pages.\footnote{Les nombres donnés pour XII à
XIX (57, 15, 24, 38, 44, 210, 200, 51) font 639 ; la page écrit « $\#\,640$ »,
le signe valant « environ ». Les titres des fascicules sont des lectures
douteuses. Les trois tomes publiés en 1972-1973 (Lecture Notes 269, 270, 305)
coupent autrement, après IV et après VIII ; ce rapprochement est le nôtre.}

La troisième donne des instructions pour un stencil : insérer le texte
manquant comme une « p. 83 bis », corriger à la page 85 des
$\mathbb{L}g^{*}$ en $\mathbb{L}g'^{*}$ et inversement --- puis tout est
barré : « Le texte oublié a été reporté ici \emph{par erreur} !!! ». Une
addition au crayon, $210 + 76 + 250 = 536$, est juste.

\pagerange{21}{27}
\section*{SGA 4 : notes exposé par exposé (pages 21 à 27)}

\subsection*{Exposé VI}

Une liste rapide (page 21) : retoucher l'introduction ; au § 4, remplacer la
notation $E_{/\varphi}$ par $\varphi\backslash E$ ; changer l'exercice sur les
topos finis ; en VI 1.25, remplacer une condition d'existence de sommes par
« il existe des limites inductives finies » ; ajouter un exercice sur les
limites inductives de topos.\footnote{Plusieurs mots de cette page sont
illisibles ; on n'a retenu que les points lisibles.}

\pagerange{22}{22}
\subsection*{Un exercice inédit : quand $\mathrm{Ind}(C)$ est-il un topos ?}

La page 22, cotée au crayon « SGA 4 VI 2 », est un exercice dactylographié qui
n'est pas dans l'exposé VI publié.\footnote{La transcription l'a vérifié sur
la réédition d'Orgogozo : VI n'y contient aucun exercice sur
$\mathrm{Ind}(C)$ comme topos, ni le mot « enveloppe de Karoubi ».}
Le voici en clair.

\emph{Exercice. Soient $C$ une petite catégorie et $C'$ une enveloppe de
Karoubi de $C$ (SGA 4 IV 7.5 b). Les conditions suivantes sont
équivalentes :}

(i) \emph{$\mathrm{Ind}(C)$ est un $\mathcal{U}$-topos ;}

(i bis) \emph{le foncteur canonique $\mathrm{Ind}(C) \to \widehat{C}$ admet un
adjoint à gauche exact à gauche ;}\footnote{La page dit « un adjoint à
droite qui est exact à gauche ». Lu à la lettre, c'est une condition vide
quant à l'exactitude : un adjoint à droite commute toujours aux limites
projectives. Ce que demande le critère de Giraud --- et ce qui équivaut à (i)
pour une sous-catégorie pleine de $\widehat{C}$ stable par limites inductives
filtrantes --- est que $\mathrm{Ind}(C)$ soit une sous-catégorie
\emph{réflexive} de $\widehat{C}$ à réflecteur exact à gauche, c'est-à-dire
que l'inclusion ait un adjoint à gauche exact à gauche (le « faisceau
associé »). Nous corrigeons en ce sens.}

(i ter) \emph{il existe une topologie sur $C$ pour laquelle les faisceaux sont
exactement les préfaisceaux ind-représentables ;}

(ii) \emph{dans $C'$ les limites inductives finies existent, le foncteur
canonique de $C'$ dans la catégorie des faisceaux y commute, pour la
topologie dont les familles couvrantes sont celles qui sont raffinées par une
famille finie épimorphique (ou épimorphique effective universelle) ;}

(iii) \emph{lorsque $C$ a des limites projectives finies, de sorte qu'on peut
prendre $C' = C$ : $C$ a des limites inductives finies et satisfait la
condition ST) de I 8.8 f) ;}

(iv) \emph{il existe un topos $E$ engendré par sa sous-catégorie pleine
$E_{PF}$ et une équivalence $C' \simeq E_{PF}$.}

\emph{De plus, la topologie de (i ter) est alors celle de (ii), le topos $E$
de (iv) est équivalent à $\mathrm{Ind}(C)$, et ce topos est parfait si et
seulement si les limites projectives finies sont représentables dans $C'$.}

Le sens moderne est clair. $\mathrm{Ind}(C)$ ne dépend que de l'enveloppe de
Karoubi --- d'où le rôle de $C'$, et la remarque juste qu'une catégorie à
limites projectives finies est déjà karoubienne. La condition (iv) dit que
$\mathrm{Ind}(C)$ est un topos \emph{localement de présentation finie} :
un topos engendré par ses objets de présentation finie, $C'$ étant
précisément la catégorie de ceux-ci, et l'on a alors
$E \simeq \mathrm{Ind}(E_{PF})$. La topologie de (ii) est la topologie
« finitaire », celle des topos cohérents.\footnote{Nous lisons $E_{PF}$
comme la sous-catégorie des objets de présentation finie ; la parenthèse qui
devait donner le renvoi est restée vide, et la condition ST) de I 8.8 f) n'est
pas rappelée sur la page. Les termes « localement de présentation finie »
(Gabriel--Ulmer, 1971) et « topologie finitaire » sont postérieurs ou
extérieurs à la page. L'énoncé est donné comme un exercice, sans
démonstration ; nous n'en avons vérifié en détail que (i) $\Leftrightarrow$
(i bis) corrigé et (i) $\Leftrightarrow$ (iv), et « parfait » est pris au
sens de VI 2.9.1, que la page ne rappelle pas.}

\pagerange{23}{24}
\subsection*{Exposés VII et IX : la cohérence du topos étale}

Page 23, deux ajouts à l'exposé VII :

\emph{5.6. (a) Le topos $\widetilde{X}_{\text{ét}}$ est localement cohérent ;
il est cohérent si et seulement si le schéma $X$ est quasi-compact et
quasi-séparé.}\footnote{La page ajoute « (resp. quasi-séparé) » après
« quasi-compact quasi-séparé », sans l'énoncé parallèle qui le justifierait ;
nous l'omettons. Les « (réf) » ajoutés au-dessus de la ligne attendent des
renvois.}

\emph{(b) Si $(X_i)_{i \in I}$ est un système projectif filtrant de schémas à
morphismes de transition affines et $X = \varprojlim X_i$ (EGA IV 8), alors
$\widetilde{X}_{\text{ét}}$ est la limite projective, au sens des topos, des
$\widetilde{X_i}_{\text{ét}}$.}\footnote{L'indice de la limite est lu « top »
sous réserve.}

\emph{Corollaire 5.6.1. Un objet $X'$ de $\mathrm{Et}(X)$ est quasi-compact
(resp. cohérent) dans $\widetilde{X}_{\text{ét}}$ si et seulement si le schéma
$X'$ est quasi-compact (resp. quasi-compact et quasi-séparé).}\footnote{La
page est elliptique (« est (quasi-compact) \ldots\ ssi $X'$ est q.c. (resp.
coh.) ») ; nous restituons l'énoncé parallèle. Elle signale aussi « VII 2.1
c) (p. 9) » comme erroné, sans dire en quoi.}

Ces énoncés servent aussitôt. La page 24 est la note annoncée pour la page 8
de l'exposé IX : elle prouve que la notion de faisceau d'ensembles
constructible de IX 2.3 coïncide avec la notion de finitude de VI 1.9.3, et
du même coup que $\widetilde{X}_{\text{ét}}$ satisfait aux conditions
équivalentes de VI 1.25. L'argument, une fois la page débarrassée d'un
premier essai qu'il a barré, est le suivant.\footnote{Un premier état d'une
quinzaine de lignes est annulé de longs traits obliques ; il identifiait la
seconde catégorie aux objets cohérents et invoquait VI 1.24. Nous ne lisons
que la version maintenue, dont l'ordre des premiers mots est lui-même
incertain.}
La question est locale ; on suppose $X$ affine, donc cohérent. Soient $C$ et
$C'$ les sous-catégories strictement pleines de $\widetilde{X}_{\text{ét}}$
formées des faisceaux constructibles au sens de IX 2.3 et au sens de VI 1.9.3 ;
il s'agit de montrer $C = C'$. Par VI 2.9 b), il suffit que $C$ soit stable
par limites projectives et inductives finies, génératrice, et formée d'objets
quasi-compacts ; les deux premiers points sont IX 2.6 (i) et 2.7.2. Reste la
quasi-compacité : par IX 2.7, tout $F \in C$ est quotient d'un $G$
représentable par un schéma quasi-compact et quasi-séparé ; $G$ est
quasi-compact dans le topos par 5.6.1 ci-dessus, et un quotient d'un objet
quasi-compact l'est encore (VI 1.3).

Il ajoute : « Bien entendu, il y aurait eu lieu de refondre les §§ 1, 2 du
présent exposé, en les rattachant aux notions générales de finitude de VI.
Nous laissons ce soin de rédaction aux générations futures. »

\pagerange{25}{27}
\subsection*{Les autres exposés}

Une page au crayon (25) note : pour IV, les modules plats sur un anneau non
commutatif et $f_{!}$ pour des foncteurs en groupes ; pour XVII 4.2.3, un
renvoi à V 3.7 peut-être changé ; pour XVIII, un théorème de finitude pour les
faisceaux d'anneaux de torsion quelconques ; une préface à la seconde édition
qui signale les autres contributions de Deligne (intégration et trace des
torseurs) ; « Que faire avec C.D. de Verdier : inclure ? » ; et pour IV § 8,
des exercices importants : topos locaux, pseudo-ponctuels,
$\widehat{I}$ pour $I$ ordonné.\footnote{« C.D. » s'entend vraisemblablement
des catégories dérivées de Verdier ; lecture de notre fait. Plusieurs mots de
la page sont incertains.}
Une demi-feuille (27) : compléter XVI 1 « comme dans Giraud ». Entre les
deux, un feuillet tapé de SGA 7 VII (les biextensions, 2.1.1 à 2.2.3), texte
publié, a servi de brouillon et porte quelques retouches à l'encre (page
26).

\pagerange{29}{30}
\section*{Une lettre sur « 2.1.1 » : le changement de base lisse pour les champs (pages 29 et 30)}

Une lettre dactylographiée, sans date, sans formule d'appel ni signature,
s'adresse à quelqu'un qu'il tutoie et dont le texte, divisé en chapitres,
contient un énoncé « 2.1.1 » sur les champs.\footnote{Le destinataire n'est pas
nommé sur les feuillets et la transcription ne le devine pas ; nous non plus.
La lettre n'a pas de date : la mention d'une « réédition de SGA 4 » la place
vers 1970-1971, ce qui n'est pas établi. Au crayon en tête : « SGA 4 XVI ».}
Sur le fond, dit-il, 2.1.1 se ramène au cas où $X$ est affine, et c'est alors
une conséquence du théorème de changement de base de SGA 4 XVI 1.1 ; mais
utiliser là toute la force du changement de base propre est « abusif » :
interprété comme un énoncé d'acyclicité globale, 2.1.1 se ramène au cas
particulier 2.1.5, puis au cas d'un corps de degré de transcendance 1, et par
là au changement de base pour le $H^1$ premier à $p$ d'une courbe propre et
lisse. Il propose de l'énoncer en toutes lettres dans la réédition de SGA 4,
ce qui dispensera le correspondant de le prouver ; sur la rédaction, il lui
reproche d'avoir « tout présent[é] à l'envers », demande de regrouper 2.1.1,
2.1.7 et 2.1.5, de remplacer « pro-lisse resp. lisse » par « localement
$L$-1-asphérique », et de reporter en IV ou V les énoncés 2.12, 2.1.3, 2.1.4,
2.1.7.1, qui « n'ont rien à voir avec les schémas ».

Le texte joint (page 30) est l'énoncé qu'il propose. Les hypothèses que la
page laisse implicites sont rétablies.\footnote{Les deux renvois à SGA 4 XV
qui donnent les exemples (« un morphisme lisse », « une limite projective
filtrante de $S$-schémas lisses à transitions affines ») ne valent que si $L$
est formé de nombres premiers inversibles sur $S$ : pour $p$ égal à une
caractéristique résiduelle, le changement de base lisse est faux (le $H^1$
d'Artin--Schreier de la droite affine). De même, l'extension
$k \subset k'$ du corollaire n'est universellement localement
$L$-1-asphérique que pour $L$ premier à la caractéristique de $k$. La page ne
le dit pas ; nous le supposons.}

\emph{Théorème 2.1.1. Soit $L$ un ensemble de nombres premiers inversibles sur
$S$, et $f : S' \to S$ un morphisme universellement localement
$L$-1-asphérique (SGA 4 XV 1.11) --- par exemple un morphisme lisse (XV 2.1),
ou un morphisme tel que $S'$ soit limite projective d'un système projectif
filtrant de $S$-schémas lisses à morphismes de transition affines (XV 1.13
(ii)). Soient $g : X \to S$ un morphisme de schémas, le carré cartésien}

\begin{tikzcd}
X \arrow[d, "g"'] & X' \arrow[l, "f'"'] \arrow[d, "g'"] \\
S & S' \arrow[l, "f"]
\end{tikzcd}

\emph{$G$ un ind-$L$-champ sur $X$ et $G' = f'^{*}G$ son image inverse sur
$X'$.}

\emph{a) Si $f$ est un morphisme local de schémas strictement locaux, le
foncteur $G(X) \to G'(X')$ est une équivalence de catégories.}

\emph{b) Si $g$ est cohérent (quasi-compact et quasi-séparé), le morphisme de
changement de base $f^{*}g_{*}G \to g'_{*}f'^{*}G$ est une équivalence de
champs sur $S'$.}

\emph{Corollaire 2.1.2. Soient $k$ un corps séparablement clos, $k'$ une
extension séparablement close de $k$, $L$ un ensemble de nombres premiers
distincts de la caractéristique, $X$ un $k$-schéma et $X' = X \otimes_k k'$.
Pour tout ind-$L$-champ $G$ sur $X$, d'image inverse $G'$ sur $X'$, le
foncteur $G(X) \to G'(X')$ est une équivalence de catégories} --- car
$\mathrm{Spec}(k') \to \mathrm{Spec}(k)$ est universellement localement
$L$-1-asphérique (renvoi à SGA 4 XVI laissé en blanc).

La démonstration est esquissée : a) par SGA 4 XVI, « nouvelle édition ! », et
V ; b) en se ramenant à a) par les lemmes 2.1.7.2 et 2.1.7.4 renumérotés.
« En tout, ça devrait faire 3 pages au Chap. VII au maximum. »

C'est le théorème de changement de base lisse en degrés $\le 1$, pour des
coefficients non commutatifs : sur un champ en groupoïdes, l'énoncé couvre à
la fois les sections ($H^0$) et les torseurs ($H^1$ non abélien au sens de
Giraud).\footnote{« Ind-$L$-champ » est le terme de la page ; nous l'entendons
comme un champ dont les faisceaux d'automorphismes sont des
ind-$L$-faisceaux, ce que la ligne tapée biffée après « champ » disait
explicitement.}

\pagerange{32}{35}
\section*{Les réponses d'Artin, et d'autres notes (pages 32 à 35)}

Une lettre d'Artin en anglais (pages 33-34), sans date ni signature sur les
feuillets, répond à ses questions sur les exposés XIII, XIX, X, XII et IX de
SGA 4 ; la transcription ne la reproduit pas. D'après sa notice, Artin accepte
l'hypothèse noethérienne en XIII 3.1 si Grothendieck la juge nécessaire,
accepte des retouches dans XIX, et fournit des références : Ax pour X, son
article « Algebraic approximation of structures over complete local rings »
pour XII 5.10, et des références pour le théorème de Tsen en IX.\footnote{La
lettre n'est pas datée. L'article qu'elle cite comme à paraître a paru aux
Publications de l'IHÉS, 36 (1969), ce qui la place vers cette date ---
datation non établie.}
La page 32, de sa main, en tire la liste : compléter des références (Giraud
pour XII, les travaux récents d'Artin pour XII 5.10, le théorème de Tsen pour
IX p. 34), revoir l'hypothèse noethérienne de XIII 3.1, des pages de XIX.

La page 35 ajoute, pour VIII, qu'il faut renvoyer à IV pour les généralités
sur les foncteurs fibres ; elle note aussi, avant de le barrer, qu'un univers
est « \emph{non vide} ! ».

\pagerange{37}{37}
\section*{Cohomologie d'un topos de préfaisceaux (page 37)}

Une page au crayon, entre des notes sur SGA 4 XVIII (Deligne) et SGA 5 VII,
propose d'ajouter à l'exposé V quelques mots sur les foncteurs dérivés de la
limite projective.

a) Sur un topos de préfaisceaux $\widehat{C}$, la cohomologie est celle de la
catégorie : pour un préfaisceau abélien $F$,
\[
H^n(\widehat{C}, F) \simeq \varprojlim\nolimits^{(n)}_{C^{\circ}} F ,
\]
les foncteurs dérivés de la limite projective sur $C^{\circ}$.\footnote{La
page écrit $\varprojlim^{(n)}_{C}$ ; $F$ étant contravariant sur $C$, la
limite est prise sur $C^{\circ}$. L'objet final de $\widehat{C}$ est le
préfaisceau constant, et $H^0 = \mathrm{Hom}(\mathbb{Z}, F) = \varprojlim F$ ;
l'énoncé est exact.}

b) Pour un topos général, il cherche une formule de type Čech :
\[
H^n(\widetilde{E}, F) \simeq \varinjlim_{\mathcal{U}}\ \varprojlim\nolimits^{(n)}_{\mathcal{U}}\, F ,
\]
la limite inductive portant sur les cribles couvrants $\mathcal{U}$ --- en
marge : « Si ça marche \ldots ».\footnote{Le membre de droite est en partie
illisible ; nous le complétons par $F$. Ainsi écrite, la formule est vraie
pour $n \le 1$ mais fausse en général : la limite sur les cribles est une
cohomologie de Čech, qui diffère de la cohomologie dérivée en degré $\ge 2$.
La forme exacte remplace les cribles par les \emph{hyperrecouvrements} ---
c'est le théorème de Verdier (SGA 4 V 7). La page ne dit pas que ça marche.}

\pagerange{38}{39}
\section*{SGA 4 VIII § 10 : étendues algébriques et espaces algébriques (pages 38 et 39)}

Un tapuscrit coté « SGA 4 VIII » porte un § 10 qui n'est ni dans l'exposé
publié ni dans le tirage original de l'IHÉS.\footnote{Vérification faite par
la transcription : l'exposé VIII publié et le tapuscrit IHÉS s'arrêtent au
§ 9. Les deux entrées manuscrites de la liste de terminologie (page 44) citent
pourtant « étendue algébrique VIII 10.2 » et « espace algébrique VIII 10.4 ».}
Il est annoncé comme une suite d'exercices « qui ne seront pas utilisés dans
la suite du Séminaire ». Les renvois à IV y sont pour la plupart laissés en
blanc.

\emph{10.1. Un schéma $X$ se reconstruit, à isomorphisme unique près, à partir
du topos annelé $(\widetilde{X}_{\text{ét}}, \mathcal{O}_X)$ connu à
équivalence près.} Les sous-objets de l'objet final du topos étale sont les
ouverts de $X$, d'où l'espace sous-jacent, $\mathrm{Top}(X) \simeq
\mathrm{Ouv}(\widetilde{X}_{\text{ét}})^{\sim}$ (VIII 6.1) ; on conclut par IV
4.2.3. Plus précisément, \emph{pour deux schémas $X$, $X'$, l'application
$\mathrm{Hom}(X, X') \to \mathrm{Homtoplocan}(\widetilde{X}_{\text{ét}},
\widetilde{X}'_{\text{ét}})$ est une équivalence de la catégorie discrète
$\mathrm{Hom}(X, X')$ sur la catégorie des morphismes de topos localement
annelés.}\footnote{La notation $\mathrm{Homtoplocan}$ est celle de IV 11
(« topos localement annelés »). On trouve ces énoncés, développés, chez M.
Hakim, \emph{Topos annelés et schémas relatifs} (1972) ; le rapprochement est
le nôtre, et rien sur la page n'y renvoie. La page écrit « à isomorphique
unique près ».}

\emph{10.2. Pour un topos annelé $(E, \mathcal{O}_E)$, les conditions suivantes
sont équivalentes :}

(i) \emph{il existe une famille $(X_i)$ d'objets de $E$ couvrant l'objet
final et des schémas $X'_i$ tels que chaque topos annelé induit $E_{/X_i}$
soit équivalent à $\widetilde{X'_i}_{\text{ét}}$ ;}

(ii) \emph{de même, avec une famille réduite à un seul objet ;}

(iii) \emph{il existe un schéma $X_0$ et une prérelation d'équivalence
$X_1 \rightrightarrows X_0$ dans la catégorie des schémas, à projection
$p_1 : X_1 \to X_0$ étale, telle que $(E, \mathcal{O}_E)$ soit équivalent au
topos annelé obtenu par recollement (descente, IV 9) du diagramme
semi-simplicial des topos étales annelés de}
\[
X_2 \mathrel{\substack{\longrightarrow\\[-2pt] \longrightarrow\\[-2pt] \longrightarrow}} X_1 \rightrightarrows X_0 ,
\qquad X_2 = (X_1, p_2) \times_{X_0} (X_1, p_1) .
\]

\emph{On dit alors que $(E, \mathcal{O}_E)$ est une étendue algébrique} ---
l'analogue des étendues topologiques de IV 9.\footnote{« Semi-simplicial »
traduit un « ½-simplicial » ajouté à l'encre et lu sous réserve. La
symétrie de la prérelation entraîne que $p_2$ est étale aussi.}

\emph{10.3. Une étendue algébrique a suffisamment de points ; elle est
localement annelée, et pour deux étendues algébriques $E$, $E'$, la catégorie
$\mathrm{Homtoplocan}(E, E')$ est un groupoïde.}

\emph{10.4. Pour une étendue algébrique $E$, les conditions suivantes sont
équivalentes :} (i) \emph{pour toute étendue algébrique $E'$,
$\mathrm{Homtoplocan}(E', E)$ est rigide, c'est-à-dire équivalente à une
catégorie discrète ;} (ii) \emph{le groupe des automorphismes de tout point
géométrique de $E$ est trivial ;} (iii) \emph{dans une présentation comme en
10.2 (iii), $(p_1, p_2) : X_1 \to X_0 \times X_0$ est un monomorphisme,
c'est-à-dire que $X_1$ est une relation d'équivalence étale.}

\emph{On dit alors que $E$ est rigide, ou que c'est un espace algébrique} ---
avec une note de sa main : « La terminologie est due à M. ARTIN, qui a mis en
évidence l'importance de cette notion en Géométrie Algébrique. » Le numéro
10.5 est tapé, et rien ne suit.

Ces pages disent, dans le langage des topos, ce que l'on dit aujourd'hui avec
des champs. Une étendue algébrique est la donnée d'un groupoïde étale de
schémas vu à travers son topos étale annelé --- ce qu'on appelle un champ de
Deligne--Mumford ; 10.3 est le fait que les morphismes de tels champs forment
des groupoïdes, et 10.4 le critère selon lequel un champ de Deligne--Mumford
à inertie triviale est un espace algébrique, quotient d'un schéma par une
relation d'équivalence étale.\footnote{Ces identifications sont les nôtres ;
elles tiennent à équivalence de topos annelés près, ce qui n'est pas tout à
fait la même chose qu'une équivalence de champs. Le tapuscrit n'est pas daté ;
l'article de Deligne et Mumford et la définition des espaces algébriques par
Artin datent de 1969, et rien sur la page ne cite le premier. Le mot
« étendue » vient de IV 9 ; dans la théorie des topos il désigne encore un
topos localement localique.}

\pagerange{40}{44}
\section*{La terminologie de SGA 4 (pages 40 à 44)}

Une liste dactylographiée, bilingue, donne pour chaque terme de SGA 4 (de VII
1.1 à XIX 6.7) son équivalent anglais et son renvoi : site et topos étale,
point géométrique, foncteur fibre, anneau hensélien et strictement local,
spécialisation et générisation, suite spectrale de Hochschild--Serre, faisceaux
de torsion et constructibles, Kummer, Hilbert 90, Artin--Schreier,
$\ell$-dimension cohomologique, fibration élémentaire, changement de base
propre et lisse, lemme de Chow, théorème de finitude, Lefschetz affine,
morphismes universellement localement $n$-acycliques et 1-asphériques,
pureté cohomologique relative et absolue, lemme d'Abhyankar, acyclicité locale,
changement de base par morphisme régulier. En tête, de sa main : « Attention,
faire \emph{deux} listes disjointes ! » --- les notations, encore mêlées aux
termes, sont entourées au crayon pour une seconde liste. Une coupure « Vol. 2 »
passe devant l'exposé XII ; deux entrées manuscrites terminent la liste :
\emph{étendue algébrique} (« algebraic spread ») et \emph{espace algébrique},
renvoyant au § 10 de VIII qu'on vient de lire.\footnote{La page 41, un premier
état abandonné des cinq premières lignes, ne porte rien de sa main. Les
entrées « théorème de finitude » ont été regroupées au crayon en une seule,
avec les renvois XIV 1.1, XVII, XVI 5.1, 5.2, XIX 5.1.}

\pagerange{46}{49}
\section*{Topos finis : un dictionnaire (pages 46 à 49)}

Au verso de deux feuillets tapés de SGA 4 et de SGA 7, une suite de notes\footnote{Le recto de la page 46 est SGA 4 IV 9.1.13, « Questions ouvertes »
sur les sous-topos ; celui de la page 48 un feuillet du second exposé sur les
biextensions de SGA 7. Tous deux sont publiés et barrés.}
compare quatre sortes d'objets finis : un espace topologique fini sobre $X$,
un ensemble ordonné fini $I$, un treillis fini $\sigma$ (les ouverts), et un
topos $E$. Le tiers supérieur de la page 46, barré, contient un premier essai ;
le reste est un « Dictionnaire » encadré, dont les flèches se lisent ainsi :

\begin{itemize}
\item $X \mapsto I$ : l'ensemble des points ordonné par générisation ;
$I \mapsto X$ : $I$ muni de la topologie dont les ouverts sont les cribles ;
\item $X \mapsto \sigma = \mathrm{Ouv}(X)$ ; $I \mapsto \sigma =
\mathrm{Crib}(I)$, et en sens inverse un foncteur $\sigma \mapsto$ (ses
« fibres », lecture incertaine) ;
\item $E \mapsto \mathrm{Ouv}(E)$ ; $\sigma \mapsto \sigma^{\sim}$, le topos
des faisceaux sur le treillis ;
\item $X \mapsto \mathrm{Top}(X)$ ; $E \mapsto \mathrm{Point}(E)$ ;
\item $I \mapsto \widehat{I}$ ; $E \mapsto \underline{\mathrm{Point}}(E)$, la
catégorie des points.
\end{itemize}

Ce sont les dualités classiques du cas fini : les espaces finis sobres
(c'est-à-dire $T_0$) correspondent aux ensembles ordonnés finis ; un treillis
distributif fini est le treillis des cribles de l'ensemble ordonné de ses
éléments premiers (dualité de Birkhoff) ; et tous ces objets ont le même
topos.\footnote{Le nom de Birkhoff et le mot « $T_0$ » sont les nôtres. Que
le foncteur inverse $\sigma \mapsto X$ soit celui des filtres premiers (les
points d'un treillis) est une lecture de notre fait : la page porte un
« $\underline{\mathrm{Fib}}(\sigma)^{\circ}$ » douteux. Une ligne isolée
définit un « schéma simplicial fini » comme un « ensemble avec ensemble
filtrant de parties » ; la définition usuelle d'un complexe simplicial
abstrait demande une famille de parties finies stable par passage aux
parties, et nous ne savons pas ce que « filtrant » visait.}
Chaque type a une opération « opposée » --- topologie opposée, ordre opposé ---,
et la page demande, entouré : « ? $E \mapsto E^{\mathrm{opp}}$ ? ». En
dessous, $\widehat{C}$ et $C^{\vee} = \widehat{C^{\circ}}$, « \emph{topos
opposé} ? ».\footnote{Une remarque de notre fait sur cette question. Pour les
topos de préfaisceaux, l'opération est bien définie à équivalence près :
$\widehat{C} \simeq \widehat{D}$ équivaut à l'équivalence des enveloppes de
Karoubi de $C$ et $D$, et l'enveloppe de Karoubi de $C^{\circ}$ est l'opposée
de celle de $C$. Elle ne s'étend pas à tous les topos, faute d'un site
privilégié, et n'est pas fonctorielle en les morphismes de topos.}

La page 48 précise le lien entre $I$ et son topos. Pour un ensemble ordonné
$I$, un \emph{crible} est une partie $U$ telle que $x \in U$ et $y \le x$
entraînent $y \in U$. Les cribles sont les ouverts d'une topologie sur $I$ ;
ce sont aussi les sous-objets de l'objet final de $\widehat{I}$, de sorte que
\[
\mathrm{Ouv}(I) = \mathrm{Crib}(I) \simeq \mathrm{Ouv}(\widehat{I}) ,
\qquad\text{d'où}\qquad
\widehat{I} \simeq \mathrm{Ouv}(I)^{\sim} = \mathrm{Top}(I) .
\]
Pour cette topologie, $y \le x$ signifie que $y$ est une générisation de $x$.
\footnote{Le « donc » de la page suppose que $\widehat{I}$ est engendré par
les sous-objets de son objet final ; c'est vrai parce que $I$ est ordonné (les
représentables $h_x \to 1$ sont des monomorphismes). Cela vaut pour tout
ensemble ordonné, fini ou non. La page porte aussi un « On a $I \simeq
\mathrm{Top}(X)$ » qui ne peut pas se lire tel quel ; nous l'omettons.}

Suit une caractérisation des topos de cette forme, en quatre lignes reliées
par des doubles flèches. Ce qui est juste est ceci : \emph{pour un topos $E$,
les conditions suivantes sont équivalentes : $E$ est engendré par des
sous-objets de l'objet final connexes, non vides et projectifs ; $E$ est
équivalent à $\widehat{I}$ pour un ensemble ordonné $I$. Pour $E =
\widehat{C}$, cela revient à dire que $C$ est un préordre (deux flèches de
même source et de même but sont égales) ; pour $E = \mathrm{Top}(X)$ avec $X$
sobre, que l'ensemble des générisations de tout point $x$ est un voisinage de
$x$} --- c'est-à-dire que $X$ est un espace d'Alexandrov.\footnote{Deux
écarts avec la page. Pour $\widehat{C}$, elle écrit « tous les morphismes de
$C$ sont des monomorphismes » ; c'est insuffisant : si $C$ est un groupe $G$,
toutes ses flèches sont inversibles mais $\widehat{C}$, le topos des
$G$-ensembles, n'a que deux sous-objets de l'objet final et n'est pas engendré
par eux dès que $G \neq 1$. Ensuite, la première des quatre lignes, « $E$ a
suffisamment de [\ldots], et est engendré par ses ouverts », a un mot
illisible que nous ne restituons pas ; si c'était « points », l'équivalence
avec la ligne suivante serait fausse (le topos des faisceaux sur
$\mathbb{R}$). Le nom d'Alexandrov est le nôtre.}

\pagerange{51}{54}
\section*{SGA 5 (pages 51 à 54)}

Des notes pour SGA 5 (pages 51-52) : en VII, le théorème de Lefschetz faible
sans supposer $X$ et $Y$ lisses, seulement $X - Y$ ; en VI, une définition
raisonnable des $\mathbb{Q}_\ell$-faisceaux en termes de champ associé, et un
nouveau numéro reprenant en langage de $\mathbb{Z}_\ell$- et
$\mathbb{Q}_\ell$-faisceaux les notions de SGA 4 ; en IV, faire dans le cas
relativement lisse toutes les constructions essentielles ; en V, signaler
l'emploi de coefficients trop particuliers ; en XIII et XIV, commenter le cas
cristallin, les formules de congruence mod $p$, et les variantes
cristallines à propos d'un théorème de Katz.\footnote{« VI » (deux fois),
« trop » et « une note » sont des lectures incertaines.}
Et, pour I § 5 : si $i : y \to X$ est l'inclusion du point générique et $F_y$
un Module localement libre sur $y$, donner
\[
D(i_{*}F_y) = i_{*}D(F_y) .
\]
\footnote{Telle quelle, la formule demande des hypothèses que la page ne donne
pas. Pour l'image directe dérivée, la règle générale est
$D \circ Rf_{!} = Rf_{*} \circ D$, et $i_{!} \neq Ri_{*}$. Ce qui est vrai,
et vraisemblablement visé, est l'énoncé pour une courbe $X$ régulière et
l'image directe \emph{non dérivée}, avec les décalages et torsions de la
dualité : c'est l'autodualité du prolongement intermédiaire $j_{!*}$, que
la littérature postérieure (Beilinson--Bernstein--Deligne) énonce en toute
dimension. Rapprochement et hypothèses sont de notre fait.}

Une lettre de R. Hartshorne, datée « Aug 4, 1969 » (pages 53-54), n'est pas
reproduite. D'après la notice de la transcription, il y demande à quel
séminaire appartiennent les notes de Jouanolou « Systèmes projectifs
$J$-adiques », exprime un doute sur la démonstration du « théorème de Shih »
qu'elles contiennent, et explique pourquoi une hypothèse artinienne y serait
inutile. Grothendieck a écrit en tête, au crayon : « Faire copie pour
Jouanolou ».

\pagerange{56}{71}
\section*{SGA 6 (pages 56 à 71)}

\subsection*{Notes}

Les notes pour SGA 6 sont courtes (pages 56 à 63) : donner en VI
l'interprétation des « suites multiplicatives de polynômes » de Hirzebruch
par une opération $f * g$ sur les séries formelles ; mentionner Serre et
Jouanolou sur la page de garde ; signaler SGA 2 V 3.5 en II 2.2.6 ; en VII,
écrire la formule d'auto-intersection pour $\mathrm{Gr}(X)$ et
$\mathrm{Gr}(Y)$ après 4.6 ; donner la structure additive explicite de
$\mathrm{Gr}(\widetilde{X}) \otimes \mathbb{Q}$ pour un schéma
éclaté ;\footnote{La page n'écrit pas de formule. Celle qu'on connaît
aujourd'hui, pour l'éclatement $\widetilde{X}$ de $X$ le long d'un
sous-schéma régulier $Y$ de codimension $c$, est
$\mathrm{Gr}(\widetilde{X})_{\mathbb{Q}} \simeq \mathrm{Gr}(X)_{\mathbb{Q}}
\oplus \bigoplus_{1 \le k \le c-1} \mathrm{Gr}(Y)_{\mathbb{Q}}$, au décalage
de graduation près ; le rapprochement est le nôtre.}
en X, tenir compte des conventions d'Illusie et faire relire par Berthelot ;
et en III 1, ajouter une note : \emph{un morphisme plat localement de
présentation finie est parfait, a fortiori pseudo-cohérent.}\footnote{Énoncé
exact : un morphisme parfait est par définition pseudo-cohérent et de
Tor-dimension finie, et un morphisme plat localement de présentation finie
l'est. Le « 6 » de « SGA 6 » est cerclé, peut-être récrit.}
La page 63 est une liste de pages d'un tapuscrit avec un numéro d'exposé (IV
ou IX) dont l'objet n'est pas dit.

Une bande déchirée (page 58) commence une construction sans contexte :
posant $\bar{\imath}\,\bar{\jmath} = \tau$, définir pour tout objet $E$ d'une
catégorie $C$ une action $\tau(E) \times E \to E$ qui fasse de $E$ un objet
formellement principal homogène sous l'objet-groupe $\tau(E)$ ; la suite est
emportée.\footnote{Lecture incertaine dans son ensemble ; la fin de la bande
manque.}

\pagerange{65}{71}
\subsection*{Commentaires d'un lecteur sur l'exposé X}

Sept feuillets dactylographiés (pages 65 à 71), « Commentaires Exp. X » et
« Commentaires Appendice », sont les remarques ligne à ligne d'un lecteur
non identifié sur la frappe d'un exposé X, selon toute apparence celui de
SGA 6, et de son appendice. Ce n'est pas un texte de Grothendieck ; la
transcription ne le reproduit pas, et rien n'y est de sa main.

\pagerange{73}{82}
\section*{SGA 7 (pages 73 à 82)}

Sous la chemise « SGA 7 » vient d'abord une lettre qui concerne SGA 6 : à
Ferrand, le 3.11.1970 (page 73), il demande s'il peut compter sur l'exposé
SGA 6 XI « d'ici fin décembre » ; sinon, il publiera SGA 6 sans lui, et
souhaite que Ferrand en fasse au moins un article d'exposition, « après le
travail que tu t'es tapé », pour pouvoir y renvoyer dans l'introduction.

Une lettre de P. Deligne, « Bures sur Yvette, le 21 octobre 1968 » (pages
74-75), n'est pas reproduite. D'après la notice de la transcription, elle
montre, sans « savants recollages de complexes cotangents », qu'un espace
analytique complexe compact $X$ est un schéma si et seulement si
$X_{\mathrm{red}}$ en est un, à l'aide du complexe cotangent et de
GAGA.\footnote{« Être un schéma » s'entend ici, selon nous, comme « être
l'analytifié d'un schéma » ; la notice ne le précise pas.}
Il a noté en tête : « Joindre SGA 7 VI App. II ? », et sa page 76 prévoit en
effet pour VI un « appendice : la lettre Deligne 2 pages ».\footnote{Le
premier mot, « Joindre », est une lecture incertaine.}

Ses notes pour SGA 7 (pages 76 à 82) prévoient encore : un appendice à XXI
sur la torsion en cohomologie cristalline (« de Katz »), en XXII des
commentaires sur la cohomologie cristalline et un renvoi à la démonstration
par Deligne du théorème de congruence de Katz généralisé, peut-être un exposé
XXIII ; dans l'appendice de VI, harmoniser la notation $L\Omega_{X/k}$ du
complexe cotangent relatif et mettre une bibliographie commune ; en I, le fait
que le groupe d'inertie a un $p$-sous-groupe de Sylow infini (renvoi à Serre,
\emph{Cohomologie galoisienne}) ; en VI, traiter aussi le « local invariant
cycle problem » pour le $H^1$, et signaler la conjecture de Griffiths en
dimension quelconque.\footnote{Le théorème des cycles invariants locaux,
établi plus tard par Deligne (\emph{La conjecture de Weil II}, 1980) en
situation de géométrie algébrique ; le rapprochement est le nôtre. « Corps »
et « Coh. galois. » sont lus sous réserve ; un bloc pour VIII, barré, parlait
de « faux schémas abéliens » et de « motifs mixtes effectifs de poids 1 ».}
En VI ou XIII, vérifier qu'on y trouve qu'en déformant une singularité
quadratique ordinaire on obtient au plus autant de telles singularités, et,
en d'autres termes, que \emph{l'ensemble des points où un morphisme plat
localement de présentation finie est lisse ou a une singularité quadratique
ordinaire est ouvert}.\footnote{« déformant », « donne », « car. rés. 2 » sont
des lectures incertaines ; un mot avant « ouvert » est illisible.}
Le reste (pages 79, 80, 82) est une liste de pages et de lignes des exposés
IX et XVIII où ajouter des références (à SGA 5 IV, à XII, XIV, XVI, au livre
de Demazure--Gabriel).

\pagerange{84}{91}
\section*{EGA I : la table de la nouvelle édition (pages 84 à 91)}

Sous la chemise « Réédition EGA I », six feuillets tapés paginés 538 à 543
sont la table des matières de la nouvelle édition d'EGA I (Springer,
\emph{Grundlehren} 166, 1971), et deux feuillets non paginés appartiennent à
son index terminologique ; textes publiés, non reproduits.\footnote{L'identification
repose sur la structure de la table (nouveau chapitre 0 avec
quasi-homéomorphismes, espaces de Jacobson, espace sobre associé ; nouveau
§ 7 du chapitre I, « Ensembles constructibles dans les schémas » ; appendice
« Ultraschémas et espaces algébriques de Serre »), non sur une collation ligne
à ligne. Les feuillets sont rangés dans l'ordre 538, 541, 539, 540, 542, 543.}
Deux marques au crayon, vraisemblablement les siennes : deux lignes périmées
(« 2.4. Fonctions constructibles », « 2.5. ») biffées au chapitre 0, et un
cadre « Morphismes affines » qui demande un numéro nouveau après I 10.15.

\pagerange{92}{97}
\section*{EGA I 10.10.5 sans hypothèse noethérienne : l'échange avec Dieudonné (pages 92 à 97)}

Le cadre est celui des schémas formels : $A$ est un anneau adique, ici
complet pour la topologie $T$-adique, $A_n = A/T^{n+1}A$,
$X = \mathrm{Spf}(A)$, $X_n = \mathrm{Spec}(A_n)$. La proposition I 10.10.5
compare, pour un $\mathcal{O}_X$-Module $F$, trois conditions : a) et c) qui
portent sur $F$ lui-même, b) qui dit que $F$ est limite projective d'un
système de Modules $F_n$ sur les $X_n$ qui se recollent, chacun de
présentation finie.\footnote{La page ne rappelle pas l'énoncé de 10.10.5 ;
nous le reconstituons seulement dans la mesure où les deux lettres de lui le
font : a) et c) équivalentes entre elles, b) plus faible hors du cas
noethérien. Le détail de a) et c) n'est pas restitué ici.}

\subsection*{La lettre du 27.8.1967 : un contre-exemple}

Dieudonné a « tiqué » sur la généralisation proposée. Grothendieck admet que
b) n'implique pas a) et c) hors du cas noethérien, et le montre. Soient $k$
un corps, $B = k[X_0, X_1, \ldots]$ l'anneau de polynômes en une infinité
d'indéterminées, $A = B[[T]]$, $A_n = B[T]/(T^{n+1})$, et
\[
J_n = (X_0,\ X_1 T,\ \ldots,\ X_n T^n) \subset A_n , \qquad M_n = A_n/J_n .
\]
Les $M_n$ se recollent ($X_{n+1}T^{n+1}$ est nul dans $A_n$) et sont de
présentation finie ; mais \emph{l'idéal $J_n$ ne peut pas être engendré par
moins de $n+1$ éléments}. Si les $M_n$ provenaient d'un $A$-module de
présentation finie $M$, on aurait des présentations
$A_n^p \to A_n^q \to M_n \to 0$ avec $p$, $q$ indépendants de $n$, et
(lemme de Schanuel) $J_n$ serait engendré par au plus $p + q$ éléments.

Preuve de la minimalité. Soit $\mathfrak{m} = (X_i)_{i \ge 0} \subset B$. Si
$a_0, \ldots, a_n \in B$ vérifient
$\sum_i a_i X_i T^i \in T J_n$ dans $A_n$, réduisons modulo l'idéal engendré
par les $X_j$, $j \neq i$ : il vient, dans $k[X_i, T]/(T^{n+1})$,
$a'_i X_i T^i \in X_i T^{i+1} k[X_i, T]$, où $a'_i$ est l'image de $a_i$ dans
$k[X_i]$ ; en comparant les coefficients de $T^i$ ($i \le n$), $a'_i X_i = 0$,
donc $a'_i = 0$ et $a_i \in (X_j)_{j \neq i} \subset \mathfrak{m}$. Ainsi les
$n+1$ générateurs écrits ont des images linéairement indépendantes sur $k$
dans $J_n / (TJ_n + \mathfrak{m}J_n)$, qui est donc de dimension $n+1$ ; par
Nakayama gradué, ou simplement parce que tout système générateur de $J_n$ en
donne un de cet espace, $J_n$ demande au moins $n+1$
générateurs.\footnote{La lettre affirme plus : que les images sont
linéairement indépendantes sur $B$ dans $J_n/TJ_n$, et conclut « d'où
évidemment $a_i = 0$ ». C'est faux : $X_0 \cdot (X_1 T) = T \cdot X_1 \cdot
X_0 \in TJ_n$, relation non triviale à coefficient $a_1 = X_0$. Son argument
(réduire modulo les $X_j$, $j \neq i$) prouve exactement que $a_i$ s'annule
\emph{dans} $k[X_i]$, c'est-à-dire $a_i \in \mathfrak{m}$, et c'est ce qui suffit
pour la minimalité ; nous rédigeons ainsi. La lettre écrit aussi
« $B[T] = A[T, (X_i)]$ » là où l'on attend $k[T, (X_i)]$.}

Il suggère donc de rédiger 10.10.5 comme deux conditions équivalentes
impliquant une troisième, équivalente aux deux autres dans le cas
noethérien, et de donner le contre-exemple.

La même lettre (page 93) signale un « canular ennuyeux » dans EGA IV 17.1.6
(i) : la démonstration du caractère local de la lissité formelle, via
16.5.17, ne marche que si $\Omega^1_{X/Y}$ est de présentation finie, par
exemple si $f$ est localement de type fini. La question équivaut, dit-il, à
celle-ci : \emph{si un $A$-module $M$ est localement projectif sur
$\mathrm{Spec}(A)$ --- s'il existe des $f_i$ engendrant l'idéal unité avec
chaque $M_{f_i}$ projectif sur $A_{f_i}$ ---, $M$ est-il projectif ?} Il
l'avait posée à D. Lazard sans se rappeler la réponse, propose un erratum, et
note qu'en pratique on prouve toujours directement la lissité formelle
globale.\footnote{La réponse est oui : la projectivité descend par
fidèle platitude, et en particulier se vérifie localement pour la topologie
de Zariski (Raynaud--Gruson, \emph{Critères de platitude et de projectivité},
Inventiones 1971, seconde partie). La lissité formelle est donc bien locale.
La lettre de 1967 ne pouvait pas le savoir ; le renvoi est le nôtre.}
Elle transmet enfin des critiques de Samuel et d'Illusie sur son projet de
§ 0 pour la réédition, et s'interrompt au milieu d'une phrase, au bas du
verso.

\subsection*{La réponse de Dieudonné, et la lettre du 15.9.1967}

Dieudonné répond le 11 septembre (page 95) ; la transcription ne reproduit
pas sa lettre. D'après sa notice, il remercie du contre-exemple mais n'est
« toujours pas convaincu » que, sans condition noethérienne, a) entraîne c),
faute de présentations $A_n^p \to A_n^q \to M_n \to 0$ à $p$ et $q$ fixes.
Grothendieck a écrit en marge : « cf. EGA IV 18.3.2.1. »

Sa réponse du 15.9.1967 (pages 96-97) accorde que les objections sont
fondées, et prouve du moins que de tels $p$ et $q$ existent, par un lemme.

\emph{Lemme. Soient $A$ un anneau, $M$ un $A$-module, $f_1, \ldots, f_m \in A$
engendrant l'idéal unité, et supposons que pour tout $i$ le $A_{f_i}$-module
$M_{f_i}$ soit engendré par $m_i$ éléments. Alors $M$ est engendré par
$\sum_i m_i$ éléments.}\footnote{La machine a laissé un blanc entre « $m$ » et
« $m_i$ » ; la suite (« $mp$ éléments, où $m$ est le nombre des $f_i$ »)
impose $\sum m_i$. Le même blanc fait que la lettre emploie $n$ pour le nombre
des $f_i$ dans l'énoncé et $m$ dans la suite ; nous écrivons $m$.}

En effet, écrivons les générateurs de $M_{f_i}$ sous la forme $h^i_j /
f_i^{N}$ avec $h^i_j \in M$ ; le sous-module $N$ engendré par tous les
$h^i_j$ vérifie $N_{f_i} = M_{f_i}$ pour tout $i$, et comme les
$\mathrm{Spec}(A_{f_i})$ recouvrent $\mathrm{Spec}(A)$, $N = M$.\footnote{La
lettre écrit « l'ouvert $\mathrm{Spec}(A_{f_i})$ de $\mathrm{Spec}(A_f)$ » et
« engendrent $M$, donc $M$ » : deux lapsus de frappe.}

Application. Dans la situation de 10.10.5, $M = \varprojlim M_n$ est un
$A$-module de type fini ; on choisit un épimorphisme $u : A^q \to M$, et des
$f_1, \ldots, f_m \in A$ dont les images engendrent l'idéal unité de $A_0$
(donc de $A$, $T$ étant topologiquement nilpotent) et tels que, sur les
ouverts correspondants de $X$, $F$ soit de présentation finie, c'est-à-dire
de la forme $\widetilde{M_i}$ avec $M_i$ de présentation finie sur l'anneau
$A_{\{f_i\}}$ de l'ouvert formel correspondant.\footnote{La lettre écrit
$A_i = A_{f_i}$ ; l'anneau de l'ouvert $\mathrm{Spf}(A)_{f_i}$ est le séparé
complété $A_{\{f_i\}}$ du localisé (EGA $0_{\mathrm{I}}$ 7.6). Cela ne change
rien à l'argument, puisque $A_{\{f_i\}} \otimes_A A_n = (A_n)_{f_i}$.}
Le noyau de $A_{\{f_i\}}^q \to M_i$ est engendré par un nombre fini $p$
d'éléments (on prend le plus grand des nombres obtenus pour les $i$) ; en
tensorisant par $A_n$,\footnote{La lettre dit « tensorisant par $A_0$ » ; la
suite porte sur $A_n$.}
la suite exacte $(A_n)_{f_i}^p \to (A_n)_{f_i}^q \to (M_n)_{f_i} \to 0$ montre
que $R_n = \mathrm{Ker}(A_n^q \to M_n)$ est localement engendré par $p$
éléments, donc par $mp$ éléments d'après le lemme --- et $mp$ ne dépend pas
de $n$.

La difficulté qui reste est de choisir ces présentations \emph{compatibles}
quand $n$ varie, c'est-à-dire un $A^p \to A^q$ d'image $\mathrm{Ker}\, u$. Si
Dieudonné y parvient, 10.10.5 retrouvera trois conditions équivalentes, en
demandant dans b) une présentation finie « uniformément en $n$ » ; sinon il
faudra un contre-exemple à a) $\Rightarrow$ c), car « la question devrait être
tirée au clair ». Le dossier n'en dit pas plus : la question reste ouverte
entre eux.

En tout état de cause, il faudra donner en corollaire, \emph{sans hypothèse
noethérienne}, l'équivalence de :

\begin{itemize}
\item[a)] $F$ est localement libre de type fini ;
\item[b)] $F$ est isomorphe à la limite projective d'une suite $(F_n)$ de
$\mathcal{O}_{X_n}$-Modules localement libres de type fini qui se recollent ;
\item[c)] $F$ est isomorphe à $\widetilde{M}$, avec $M$ un $A$-module
projectif de type fini,
\end{itemize}

en procédant comme pour 10.10.5 avec EGA IV 18.3.2.1 --- un lemme qui, dit-il,
devrait venir en corollaire après $0_{\mathrm{I}}$ 7.2.9.\footnote{L'énoncé
est exact ($X = \mathrm{Spf}(A)$ affine, $A$ complet et séparé) : un module
projectif de type fini sur $A_0$ se relève, de façon unique à isomorphisme
près, en un module projectif de type fini sur chaque $A_n$, puis sur $A$, et
$M = \varprojlim M_n$.}
Quant aux modifications qu'il préconisait pour I 10.11, elles tombent si
10.10.5 ne s'arrange pas ; on peut cependant dire qu'un $F$ de présentation
finie est limite projective de $F_n$ de présentation finie qui se recollent
(sans réciproque), que $F$ est localement libre si et seulement si les $F_n$
le sont, et que le reste de 10.11 vaut sans hypothèse noethérienne, sauf deux
énoncés.\footnote{« ssi » est une frappe surchargée, lue sous réserve. Les
deux exceptions sont « 10.7.11.2 » --- ainsi sur la page, où l'on attendrait
un numéro de 10.11 --- et la partie « injectivité » de 10.11.9.}

Entre les deux lettres, une page (94) note seulement, pour I 10.12
(« Morphismes adiques de schémas formels ») : « Stabilité par
$\varprojlim$ (ou non ?) ».

\pagerange{99}{99}
\section*{EGA II (page 99)}

Une dernière page énumère pour la réédition d'EGA II : (a) transférer
EGA IV 4.8, un exposé des morphismes affines de préschémas formels, un
appendice sur les schémas formels ; (b) en II 6.5.10 et 6.5.11, ajouter des
formules, la seconde pour le cas où $X' \to X$ est génériquement décomposé, et
en 6.5.7 renvoyer à EGA IV 21 ; (c) remonter EGA IV 19.4.11, « revu et
augmenté » ; (d) un point sur 7.4.2.\footnote{Plusieurs mots sont illisibles
ou incertains : « écrit », « préschémas », « pour », « décomposé », et la
ligne (d) presque entière.}

\end{document}
